% Options for packages loaded elsewhere
\PassOptionsToPackage{unicode}{hyperref}
\PassOptionsToPackage{hyphens}{url}
\PassOptionsToPackage{dvipsnames,svgnames,x11names}{xcolor}
%
\documentclass[
  11pt,
]{article}
\usepackage{amsmath,amssymb}
\usepackage{iftex}
\ifPDFTeX
  \usepackage[OT1]{fontenc}
  \usepackage[utf8]{inputenc}
  \usepackage{textcomp} % provide euro and other symbols
\else % if luatex or xetex
  \usepackage{unicode-math} % this also loads fontspec
  \defaultfontfeatures{Scale=MatchLowercase}
  \defaultfontfeatures[\rmfamily]{Ligatures=TeX,Scale=1}
\fi
\ifPDFTeX\else
  % xetex/luatex font selection
\fi
% Use upquote if available, for straight quotes in verbatim environments
\IfFileExists{upquote.sty}{\usepackage{upquote}}{}
\IfFileExists{microtype.sty}{% use microtype if available
  \usepackage[expansion=false]{microtype}
  \UseMicrotypeSet[protrusion]{basicmath} % disable protrusion for tt fonts
}{}
\makeatletter
\@ifundefined{KOMAClassName}{% if non-KOMA class
  \IfFileExists{parskip.sty}{%
    \usepackage{parskip}
  }{% else
    \setlength{\parindent}{0pt}
    \setlength{\parskip}{6pt plus 2pt minus 1pt}}
}{% if KOMA class
  \KOMAoptions{parskip=half}}
\makeatother
\usepackage{xcolor}
\usepackage[margin=1in]{geometry}
\usepackage{longtable,booktabs,array}
\usepackage{calc} % for calculating minipage widths
% Correct order of tables after \paragraph or \subparagraph
\usepackage{etoolbox}
\makeatletter
\patchcmd\longtable{\par}{\if@noskipsec\mbox{}\fi\par}{}{}
\makeatother
% Allow footnotes in longtable head/foot
\IfFileExists{footnotehyper.sty}{\usepackage{footnotehyper}}{\usepackage{footnote}}
\makesavenoteenv{longtable}
\setlength{\emergencystretch}{3em} % prevent overfull lines
\providecommand{\tightlist}{%
  \setlength{\itemsep}{0pt}\setlength{\parskip}{0pt}}
\setcounter{secnumdepth}{-\maxdimen} % remove section numbering
\ifLuaTeX
  \usepackage{selnolig}  % disable illegal ligatures
\fi
\usepackage{bookmark}
\IfFileExists{xurl.sty}{\usepackage{xurl}}{} % add URL line breaks if available
\urlstyle{same}
\hypersetup{
  pdftitle={Exact Divisor Densities and a Finite-to-Global Certification of Integer Maximizers},
  pdfauthor={Steven Minh Phat Duong, MPH},
  colorlinks=true,
  linkcolor={Maroon},
  filecolor={Maroon},
  citecolor={Blue},
  urlcolor={Blue},
  pdfcreator={LaTeX via pandoc}}

\title{Exact Divisor Densities and a Finite-to-Global Certification of
Integer Maximizers}
\usepackage{etoolbox}
\makeatletter
\providecommand{\subtitle}[1]{% add subtitle to \maketitle
  \apptocmd{\@title}{\par {\large #1 \par}}{}{}
}
\makeatother
\subtitle{The cases n = 4 through 20}
\author{Steven Minh Phat Duong, MPH}
\date{18 November 2026 -- Preprint; availability correction}

\begin{document}
\renewcommand{\labelitemi}{$\bullet$}
\maketitle

Independent Researcher, Calgary, Alberta, Canada

\subsection{Abstract}\label{abstract}

We describe an exact-arithmetic certificate scheme for identifying the
integer endpoint that maximizes the density of integers having exactly
one divisor in an open interval \((n,m)\). Two elementary bounds reduce
an infinite endpoint range to a finite sweep, a chain of bounded
intervals and one tail inequality. We derive three prime-valuation
algorithms for the required densities and report certificate-based
candidate classifications for each fixed \(n=4,\ldots,20\). The listed
maximizing endpoints are
\(9,12,12,15,18,20,22,24,24,32,32,32,32,38,38,44,44\). The
finite-to-global criterion is proved here; the numerical classification
depends on the correctness of the supplied exact density entries. All 17
logical certificates were rechecked. The accompanying validation record
distinguishes complete Program C recomputation from comparisons with
saved Program A logs. An additional inclusion-exclusion implementation
reproduces all 17 peak values without using the prime-valuation
decomposition. The work does not claim a general formula for maximizers,
a new general existence theorem, independent human review or established
priority.

\subsection{1. Definitions and result
scope}\label{definitions-and-result-scope}

Fix integers \(n\ge1\) and \(m>n+1\). Put
\[D(n,m)=\{d\in\mathbb Z:n<d<m\},\qquad c_m(x)=|\{d\in D(n,m):d\mid x\}|.\]
For \(j=0,1\), define
\[\delta_j(n,m)=\lim_{X\to\infty}\frac{|\{1\le x\le X:c_m(x)=j\}|}{X}.\]
These events are periodic with period \(L=\operatorname{lcm}(D(n,m))\),
so the limits exist and are rational. For an empty interval use \(L=1\),
\(\delta_0=1\), \(\delta_1=0\).

The certificate-reported classification is: for each
\(n\in\{4,\ldots,20\}\), the row of Table 1 below gives the unique
integer \(m_*>n+1\) for which
\[\delta_1(n,m)<\delta_1(n,m_*)\quad\text{for every integer }m>n+1,\ m\ne m_*.\]
This conclusion depends on both the reduction below and the correctness
of the exact densities in the certificates. The logical checker verifies
implications between supplied values; it does not establish their
interpretation as densities.

For a real endpoint \(t\), \(D(n,t)=D(n,\lceil t\rceil)\). Thus the
density is constant on \((m-1,m]\). Uniqueness concerns the integer
endpoint, or equivalently the divisor set, rather than an individual
real endpoint.

\subsection{2. Bounding lemmas and complete endpoint
coverage}\label{bounding-lemmas-and-complete-endpoint-coverage}

Write \(E_j(m)=\{x:c_m(x)=j\}\). For integers \(a\le m\le b\), nested
divisor sets give \(c_a\le c_m\le c_b\) pointwise.

\textbf{Block lemma.} For \(a\le m\le b\),
\[\delta_1(n,m)\le\delta_1(n,a)+\delta_0(n,a)-\delta_0(n,b).\] Indeed,
\(c_m(x)=1\) implies \(c_a(x)\in\{0,1\}\) and \(c_b(x)\ge1\), so
\[E_1(m)\subseteq E_1(a)\cup\bigl(E_0(a)\setminus E_0(b)\bigr).\] The
two sets on the right are disjoint and \(E_0(b)\subseteq E_0(a)\).
Taking densities proves the bound, including both endpoints.

\textbf{Tail lemma.} For \(m\ge M\),
\[\delta_1(n,m)\le\delta_0(n,M)+\delta_1(n,M).\] Here \(c_m(x)=1\)
implies \(c_M(x)\in\{0,1\}\), giving
\(E_1(m)\subseteq E_0(M)\cup E_1(M)\). Take densities of these disjoint
events.

\textbf{Proposition (finite-to-global certificate criterion).} Fix \(n\)
and integers \(n+2\le m_*\le S\le M\). Suppose the exact densities at
all sweep endpoints \(n+2,\ldots,S\) are supplied, and let
\(B=\delta_1(n,m_*)>0\). Suppose (i) every other swept value is strictly
less than \(B\); (ii) there is a contiguous chain of closed integer
blocks from \(S\) to \(M\) whose block-lemma bounds are strictly less
than \(B\); and (iii) \(\delta_0(n,M)+\delta_1(n,M)<B\). Then \(m_*\)
uniquely maximizes \(\delta_1(n,m)\) over all integers \(m>n+1\).

\emph{Proof.} Every admissible endpoint belongs to the sweep, a block,
or the tail. Condition (i) handles the sweep. The block lemma and
condition (ii) handle the blocks. The tail lemma and condition (iii)
handle every endpoint at least \(M\). These sets cover all admissible
endpoints, and every endpoint except \(m_*\) has density strictly below
\(B\). For an empty block chain, \(S=M\) and the sweep and tail suffice.
\(\square\)

Each certificate supplies the exact sweep \(m=n+2,\ldots,1000\), closed
integer blocks from 1000 to a tail point \(M\), and both densities at
\(M\). Every nonpeak sweep value, every block bound and the tail bound
must be strictly below the claimed peak. Adjacent blocks share
endpoints. Sweep, blocks and tail therefore cover every admissible
integer endpoint without gaps. For \(n=4\) the block list is empty and
\(M=1000\). The largest \(M\) is 256875, for \(n=19,20\). Endpoints
\(m\le n+1\), if considered, give zero and cannot tie a positive peak.

\newpage

\subsubsection{Table 1. Certificate-reported integer
maximizers}\label{table-1.-certificate-reported-integer-maximizers}

\begin{longtable}[]{@{}rrlr@{}}
\toprule\noalign{}
\(n\) & \(m_*\) & \(\delta_1(n,m_*)\) & Tail \(M\) \\
\midrule\noalign{}
\endhead
\bottomrule\noalign{}
\endlastfoot
4 & 9 & 5/14 & 1000 \\
5 & 12 & 1109/3465 & 2875 \\
6 & 12 & 1172/3465 & 1875 \\
7 & 15 & 13532/45045 & 8875 \\
8 & 18 & 14048/45045 & 8875 \\
9 & 20 & 497888/1616615 & 8875 \\
10 & 22 & 528/1729 & 16875 \\
11 & 24 & 10954544/37182145 & 32875 \\
12 & 24 & 11091536/37182145 & 32875 \\
13 & 32 & 137993457152/501401225325 & 128875 \\
14 & 32 & 470752768/1676927175 & 64875 \\
15 & 32 & 47158954496/167133741775 & 128875 \\
16 & 32 & 1638932224/5763232475 & 64875 \\
17 & 38 & 1697006237184/6183948445675 & 128875 \\
18 & 38 & 1733299766784/6183948445675 & 128875 \\
19 & 44 & 2897060040192/10741183359335 & 256875 \\
20 & 44 & 2768951781888/10048203787765 & 256875 \\
\end{longtable}

The table is a transcription of the unchanged certificate-generated
results. The classification is conditional on the density values being
correctly computed.

\subsection{3. Common exact measure}\label{common-exact-measure}

Prime valuations have densities
\[\Pr(v_p(x)=a)=(1-1/p)p^{-a},\qquad\Pr(v_p(x)\ge a)=p^{-a}.\] Only
finitely many capped valuations matter: cap at the largest exponent
occurring in the relevant divisors. The resulting events depend on
residue classes modulo prime powers. The Chinese remainder theorem
multiplies their densities, giving independence across these primes.
Primes greater than \(m-1\) divide no member of \(D(n,m)\). All three
algorithms share this model, so their agreement is not an independent
test of it.

\subsection{4. Program A by avoidance
recursion}\label{program-a-by-avoidance-recursion}

For a finite set \(A\) of positive integers, let \(F(A)\) be the density
of integers divisible by no member of \(A\). Then \(F(\varnothing)=1\)
and \(F(A)=0\) if \(1\in A\). Removing an element divisible by another
member of \(A\) leaves \(F(A)\) unchanged.

Partitioning by the unique divisor gives
\[\delta_0(n,m)=F(D(n,m)),\qquad\delta_1(n,m)=\sum_{d\in D(n,m)}\frac1d F(A_d),\]
where \(A_d=\{e/\gcd(e,d):e\in D(n,m),e\ne d\}\). Write \(x=dy\):
\(e\mid dy\) iff \(e/\gcd(e,d)\mid y\). Multiplication by \(d\) scales
density by \(1/d\), and the events for different unique divisors are
disjoint.

For nonempty \(A\) not containing 1, let \(p\) be its largest prime
factor, \(K=\max_{a\in A}v_p(a)\), and
\[A_k=\{a/p^{v_p(a)}:a\in A,v_p(a)\le k\},\qquad B=\{a/p^{v_p(a)}:a\in A\}.\]
Conditioning on \(v_p(y)\) gives
\[F(A)=\sum_{k=0}^{K-1}\frac{p-1}{p^{k+1}}F(A_k)+\frac1{p^K}F(B).\] The
last term sums all exponents at least \(K\), whose constraints coincide.
Every recursive set is \(p\)-free, so the largest relevant prime
decreases and recursion terminates. Empty sets and sets containing 1
cover the boundary cases; memoization affects cost only.

Source anchors: \texttt{progA/delta.py}: \texttt{pfac},
\texttt{primitive}, \texttt{avoid}, \texttt{delta0}, \texttt{delta1};
\texttt{progA/sweepA.py} records the sweep. The floating-point helper
\texttt{bound\_crossing} is outside the exact certificate path and is
not used for these claims.

\subsection{5. Program B by the square-root prime
split}\label{program-b-by-the-square-root-prime-split}

Put \(h=m-1\). Small primes satisfy \(p^2\le h\); large primes satisfy
\(q\le h<q^2\). A divisor \(d\le h\) contains at most one large prime,
to exponent 1: a square or product of two such primes exceeds \(h\).

Condition on the small-prime valuations, and let \(S\) be that part of
\(x\). Let \(w_0\) count divisors of \(S\) in \([n+1,h]\) and \(w_q\)
count divisors \(s\mid S\) with \(sq\in[n+1,h]\). Then
\[c_m(x)=w_0+\sum_{\substack{q\text{ large}\\q\mid x}}w_q.\] For this
state put \(\Pi=\prod_{w_q\ge1}(1-1/q)\) and \(A_1=\{q:w_q=1\}\). If
\(w_0=0\), the conditional contributions to \((\delta_0,\delta_1)\) are
\((\Pi,\Pi\sum_{q\in A_1}1/(q-1))\). If \(w_0=1\), they are \((0,\Pi)\).
If \(w_0\ge2\), both are zero. These follow from independent
divisibility by the large primes.

For each small prime \(p\), let \(a_{\max}=\max\{a:p^a\le h\}\).
Enumerate exact exponents below \(a_{\max}\) with weights
\((p-1)/p^{a+1}\) and the tail with weight \(p^{-a_{\max}}\). Higher
exponents create no additional divisor at most \(h\). Each DFS step
generates all divisors of the partial product at most \(h\); larger
divisors can be discarded because later multiplication cannot shrink
them. Once two divisors lie in the interval, every extension still has
at least two. That state contributes to neither density, justifying
pruning; the same monotonicity in the current exponent justifies the
loop break.

Source anchors: \texttt{progB/verify.py}: \texttt{dens}, its nested
\texttt{dfs}/\texttt{leaf}, and \texttt{set\_prime\_limit}. The current
code \textbf{enforces} \(n\ge1\) and \(m-1\le\texttt{PR\_N}\). Its
default limit is 1000; \texttt{gen\_cert.py} sets it to \texttt{mmax+10}
before larger runs. These guards do not change successful density
values.

\subsection{6. Program C by primes in the open
interval}\label{program-c-by-primes-in-the-open-interval}

Let \(s(x)\) be the part of \(x\) supported on primes at most \(n\), and
\(\mathcal Q=\{Q\text{ prime}:n<Q<m\}\). Two distinct primes from
\(\mathcal Q\) already give two divisors in \(D(n,m)\), so only zero or
one can contribute.

For \(q\le n\) prime, let \(c_q\) be the least exponent with
\(q^{c_q}\ge m\). Use exact weights below \(c_q\) and tail weight
\(q^{-c_q}\) at \(c_q\). Higher exponents enter no relevant divisor. Let
\(S_j\) be the total small-prime weight for exactly \(j\) divisors in
\(D(n,m)\) and \(P_n=\prod_{Q\in\mathcal Q}(1-1/Q)\). With no interval
prime, the contributions are \(P_nS_0\) and \(P_nS_1\).

If exactly one interval prime \(Q\) divides \(x\), \(Q\) is the required
divisor. There must be no divisor of \(s(x)\) in the interval. If
\(Q^2<m\), its valuation must be exactly 1; otherwise any positive
valuation is allowed. Every prime factor \(q\) of \(s(x)\) must satisfy
\(qQ\ge m\), or \(qQ\) would be a second divisor. Conversely, these
restrictions exclude every other mixed divisor: for \(s>1\),
\(sQ^a\ge qQ\ge m\), and powers of \(Q\) alone are excluded by the
valuation restriction.

Put \(\theta=\lceil m/Q\rceil\),
\(R_\theta=\{q\le n:q\text{ prime},q\ge\theta\}\) and
\[Z_\theta=W_0(R_\theta)\prod_{\substack{q\le n\text{ prime}\\q<\theta}}(1-1/q),\]
where \(W_0(R_\theta)\) is the weight of states supported on
\(R_\theta\) having no interval divisor. Dividing the permitted
\(Q\)-valuation weight by its factor \(1-1/Q\) in \(P_n\) gives
\[\delta_0=P_nS_0,\qquad\delta_1=P_n\left(S_1+\sum_{Q\in\mathcal Q}Z_{\lceil m/Q\rceil}\alpha_Q\right),\quad
\alpha_Q=\begin{cases}1/Q&Q^2<m,\\1/(Q-1)&Q^2\ge m.\end{cases}\]
Grouping equal allowed-prime sets and using tree sums/products changes
only evaluation order. The DFS adds disjoint divisor layers as exponents
increase; the count is monotone, so it can prune at two interval
divisors (or one for \(W_0\)). Empty intervals give \((1,0)\), and
\(n=1\) has an empty small-prime set. The code enforces \(n\ge1\).

Source anchors: \texttt{progC/progC.py}: \texttt{compute}, nested
\texttt{wt}/\texttt{enum}, \texttt{tree\_prod} and
\texttt{tree\_sum\_inv}.

\subsection{7. Evidence and the trust
boundary}\label{evidence-and-the-trust-boundary}

\texttt{progB/gen\_cert.py} generated 17 compressed certificates.
\texttt{checker/check\_cert.py} uses exact fractions and no density
engine: it checks coverage, a unique swept maximum, consistency at
repeated points, and strict block/tail inequalities.
\texttt{checker/compare\_C.py} recomputes the certificate values with C;
\texttt{checker/compare\_A.py} compares against saved A logs. Default
\texttt{replay.sh} reruns C and the validators, but \textbf{does not
rerun A}. Full mode reruns all 16779 swept values with A and costs
several CPU-hours. The dated validation record specifies the replay mode
completed for this version; full mode is distinct from default replay.

The stored bundle reports period-counting checks of 197 B cases, 188 C
cases and 5166 lemma cases, with no mismatches, plus logical, negative
and failure-propagation tests. Those are recorded checks. The earlier
audit also reports A/B/C agreement on 900 small pairs; its original text
is retained in the earlier review snapshot. The earlier audit was
AI-assisted, not independent human review.

The algorithms share the valuation model, Python integer arithmetic and
\texttt{Fraction}. Different decompositions reduce some risks but do not
eliminate correlated errors. C had an early bug caught by brute-force
validation and corrected before certificate comparisons.
\texttt{legacy/} runs and original reports remain available; their
exploratory sweep/tail parameters do not describe the current
certificates. Successful computational checks do not establish human
mathematical acceptance or novelty.

\subsection{8. Relation to earlier work}\label{relation-to-earlier-work}

Cambie's \href{https://arxiv.org/abs/2501.10333}{paper on unimodularity}
addresses unimodality. His
\href{https://github.com/StijnCambie/ErdosProblems/blob/main/Erdosproblem692/n_le20.ipynb}{n\_le20
notebook} displays rounded values for \(m=n+2,\ldots,n+15\). For
\(n=4,\ldots,12\) the candidate peak locations are visible in that
window; exclusion of all later endpoints is a different claim. For
\(n=13,\ldots,20\) the candidate maximizers lie outside it.
\href{https://github.com/KitaKen1/erdos-692-part-ii-lean}{KitaKen1's
Lean development} states general existence of a maximum. That theorem is
not claimed here as new, and its Lean build was not replayed. See also
the \href{https://www.erdosproblems.com/692}{problem statement}. Novelty
remains provisional.

The proposed contribution is a reproducible finite-parameter
classification with exclusion of all later integer endpoints. It is not
the general existence theorem or a resolution of the unimodality
question. A comprehensive later-literature and repository comparison
remains outstanding, and no priority claim is made. The accessible
prior-work sources were checked on 18 November 2026; the Lean proof was
not built.

\subsection{9. Reproducibility and
limitations}\label{reproducibility-and-limitations}

Use \textbf{Python 3.11+}, standard library only, for the certificate
pipeline. \texttt{gen\_cert.py} uses
\texttt{sys.set\_int\_max\_str\_digits}; exact fractions can contain
very large integers. Readers disable the decimal-digit conversion limit
where needed. Replay scripts use a POSIX shell; integrity can be checked
with \texttt{sha256sum} or macOS \texttt{shasum\ -a\ 256}.

Historical audit records are preserved separately. The supplied
certificates, numerical results and original density algorithms are
unchanged. New validation and packaging changes are recorded explicitly.

\subsubsection{Validation of the packaged
computation}\label{validation-of-the-packaged-computation}

The accompanying package includes all 17 certificates and all 17 saved
Program A sweep logs. The logical checker verifies coverage and strict
rational inequalities; it does not recompute densities. Default replay
separately validates the density engines on small periods, recomputes
every certificate density using Program C, compares every swept value
with the saved Program A output, and tests that deliberately corrupted
copies fail. Full mode additionally regenerates Program A's complete
sweep; it must not be conflated with default replay. On 18 November
2026, the complete default replay finished with exit status zero: all 17
Program C certificate comparisons and all saved-A comparisons had zero
discrepancies, all logical and small-period checks passed, and all three
end-to-end corruption tests failed as intended. The 14 negative tests
and five logical-checker fixtures also passed. Full-mode regeneration of
Program A was not performed. Logs are supplied in \texttt{validation/}
and \texttt{RELEASE-STATUS.md}.

As an additional check, weighted inclusion-exclusion was implemented
using least-common-multiple states. For every finite divisor set \(D\),
\[\delta_1=\sum_{\varnothing\ne A\subseteq D}\frac{(-1)^{|A|+1}|A|}{\operatorname{lcm}(A)}.\]
For an integer divisible by exactly \(r\) members of \(D\), its
coefficient is \(\sum_{k=1}^{r}(-1)^{k+1}k\binom{r}{k}\), equal to one
when \(r=1\) and zero otherwise. Averaging periodic divisibility
indicators gives the formula. Inclusion-exclusion is standard and also
appears in Cambie's notebook; it is not claimed as new.

The new implementation groups subsets by their LCM. For each LCM \(L\),
it stores \(a_L=\sum_A(-1)^{|A|}\) and \(b_L=\sum_A(-1)^{|A|}|A|\).
Initially \((a_1,b_1)=(1,0)\). When a divisor \(d\) is added, exclusion
preserves the old pair; inclusion adds \((-a_L,-b_L-a_L)\) at
\(\operatorname{lcm}(L,d)\). The required density is \(-\sum_L b_L/L\).
This recurrence reproduces the 17 exact peak fractions and agrees with
54 direct finite-period counts. It independently checks peak heights,
not exclusion of every later endpoint. Like the other implementations,
it uses Python integers and rational arithmetic and was AI-assisted.

\subsubsection{Data and code
availability}\label{data-and-code-availability}

The preprint and supporting v3 archive are available in the published
Zenodo record \href{https://doi.org/10.5281/zenodo.24338898}{10.5281/zenodo.24338898},
deposited on 18 November 2026 with record version label \texttt{1.0-draft}.
The archive contains the mathematical derivation, unchanged exact
certificates, density engines, saved logs, inclusion-exclusion checkers,
replay scripts and integrity manifests. This corrected manuscript updates
availability and publication-status wording only; it does not change the
mathematical claims or numerical results. 

The record displays a CC BY 4.0 licence. This reports the repository
metadata; it does not independently establish ownership of every included
file, resolve any adapted-source obligations, or override applicable
third-party notices, including those for embedded fonts. Publication of a
preprint is not peer review.

The historical full audit archive has SHA-256

\texttt{31a4b78ded4bb9f18e3f8fb3f0470d9f39130be348ec047c70446cc42cfba777}.

The exact supporting-archive file inventory is in its manifest. The
original archive and earlier review snapshots are preserved separately.
This manuscript does not imply that a repository's latest branch was
audited.

\subsubsection{AI assistance and
responsibility}\label{ai-assistance-and-responsibility}

AI tools assisted the mathematical development, implementations,
computational runs, audits and manuscript preparation. Agreement between
AI-assisted implementations is not independent human mathematical
review. Responsibility for the mathematical claims and interpretation
rests with the author. No faculty endorsement, institutional
affiliation, peer review or formal verification of this classification
is claimed. This manuscript is a preprint, not a peer-reviewed article.

\subsection{References}\label{references}

\begin{enumerate}
\def\labelenumi{\arabic{enumi}.}
\tightlist
\item
  Stijn Cambie. \emph{Resolution of Erdős' problems about
  unimodularity}. Journal of Number Theory \textbf{280} (2026),
  271--277. \href{https://doi.org/10.1016/j.jnt.2025.08.014}{Published
  article}; \href{https://arxiv.org/abs/2501.10333}{accessible
  preprint}.
\item
  Stijn Cambie. \emph{ErdosProblems}, problem 692 notebooks.
  \href{https://github.com/StijnCambie/ErdosProblems/tree/main/Erdosproblem692}{Repository}.
\item
  KitaKen1. \emph{Erdős Problem 692, Part II: existence of a maximum},
  Lean repository.
  \href{https://github.com/KitaKen1/erdos-692-part-ii-lean}{Repository}.
\item
  \emph{Erdős Problems}, problem 692. https://www.erdosproblems.com/692
\end{enumerate}

\end{document}
